\newgeometry{left=25mm, right=35mm, top=25mm, bottom=30mm}
\chapter{Systematische Literaturrecherche}
\label{ch:bsp}
\subsubsection{Einleitung \& Forschungsfrage} 
Im Rahmen dieser Arbeit wird eine systematische Literaturrecherche (SLR) durchgeführt, um ein umfassendes Verständnis der aktuellen Forschung im Bereich der Anomalieerkennung \\ mittels generativer KI in Echtzeit zu gewinnen. Die Relevanz dieses Themas ergibt sich aus der zunehmenden Datenmenge in modernen Systemen und der Notwendigkeit, Abweichungen von erwarteten Mustern sogenannte 'Anomalien' schnell und präzise zu identifizieren. Insbesondere in dynamischen Umgebungen, in denen sich Datenmuster ständig weiterentwickeln können, stellt die Unterscheidung zwischen echten Anomalien und lediglich neuen, aber validen Datenpunkten eine signifikante Herausforderung dar.
\\
Die zentrale Forschungsfrage, die in dieser systematischen Literaturrecherche untersucht wird, lautet:
\\
Wie kann sichergestellt werden, dass Daten in Echtzeit mithilfe von generativer KI als Anomalie klassifiziert werden können, und zwar unter Berücksichtigung der Herausforderung, dass es sich hierbei wirklich um eine Anomalie handelt und nicht um einen neuen, validen Datenpunkt?
\\
Diese Frage zielt darauf ab, Methoden, Techniken und Modelle zu identifizieren, die eine \\ zuverlässige und kontextsensitive Anomalieerkennung ermöglichen. Dabei liegt ein besonderer Fokus auf dem Potenzial generativer KI-Ansätze, die adaptiv auf sich ändernde Daten reagieren können, ohne fälschlicherweise neue, legitime Datenpunkte als Anomalien zu interpretieren.
\nAbsatz
Die Durchführung dieser systematischen Literaturrecherche folgt einem präzisen und \\ nachvollziehbaren methodischen Vorgehen, um die Objektivität und Reproduzierbarkeit \\ der Ergebnisse zu gewährleisten. Die Methodik umfasst die Auswahl relevanter Datenbanken, die Definition geeigneter Suchbegriffe, Kriterien für die Ein- und Ausschluss von Publikationen sowie einen detaillierten Auswahlprozess der relevanten Literatur.
\nAbsatz
Für die Identifizierung relevanter Publikationen wird überwiegend Google Scholar als zentrale Suchmaschine genutzt. Google Scholar bietet eine breite Abdeckung wissenschaftlicher Literatur aus verschiedenen Disziplinen und ermöglicht durch seine Suchalgorithmen eine effiziente \\ Identifizierung relevanter Studien. Es werden vorzugsweise Veröffentlichungen  IEEE Xplore oder ACM Digital Library benutzt.
\\
Die Suchstrategie basiert auf einer Kombination von Schlüsselwörtern, die die Kernkonzepte der Forschungsfrage abbilden. Es werden sowohl deutsche als auch englische Begriffe verwendet, um eine möglichst umfassende Erfassung relevanter Literatur zu gewährleisten. \\\textbf{ Die Hauptschlagwörter umfassen:} 
\nAbsatz
    \textbf{Deutsche Suchbegriffe: }\\
        'Anomalieerkennung mit KI', \\
        'Anomalieerkennung in adaptiven Systemen mit KI',\\
        'Echtzeit Anomalieerkennung KI',\\
        'Generative KI Anomalie'\\
\newpage
    \textbf{Englische Suchbegriffe:}\\
        'Anomaly detection with AI',\\
        'Anomaly detection in adaptive systems with AI',\\
        'Real-time anomaly detection AI',\\
        'Generative AI anomaly',\\
        'Novel data point classification AI',\\
        'Outlier detection generative models'\\
\\
Diese Suchbegriffe werden in verschiedenen Kombinationen verwendet, um die Suchergebnisse zu optimieren und irrelevanten Publikationen vorzubeugen.
\subsubsection{Einschluss- und Ausschlusskriterien:}
\\
Um die Qualität und Relevanz der ausgewählten Literatur sicherzustellen, \\ werden präzise Einschluss- und Ausschlusskriterien definiert:
\nAbsatz
\textbf{Einschlusskriterien:}
\nAbsatz
    \textbf{Sprache:} Es werden ausschließlich Publikationen berücksichtigt, die in deutscher oder \\ englischer Sprache verfasst sind.\nAbsatz
    \textbf{Forschungsbereich:} Die Studien sollten vorzugsweise aus der naturwissenschaftlichen \\ Forschung stammen, insbesondere aus den Bereichen Informatik, Künstliche Intelligenz, Data Science, Ingenieurwissenschaften oder angrenzenden Gebieten. Dies ist wichtig, da in diesen \\ Disziplinen Konzepte und Daten in der Regel präziser definiert und reproduzierbar sind, \\ im Gegensatz zu den oft interpretativeren Ansätzen der Geisteswissenschaften.
    \nAbsatz
    \textbf{Inhaltsrelevanz:} Die Publikationen müssen einen direkten Bezug zur Anomalieerkennung, generativer KI, Echtzeit-Datenverarbeitung und der Unterscheidung von Anomalien und neuen validen Datenpunkten aufweisen.
    \nAbsatz
   \textbf{Publikationsart:} Bevorzugt werden wissenschaftliche Artikel (Peer-Reviewed), \\ Konferenzbeiträge, Dissertationen und anerkannte Fachbücher.
\nAbsatz
\textbf{Ausschlusskriterien:}
\nAbsatz
    \textbf{Sprache:} Publikationen in anderen Sprachen als Deutsch oder Englisch, \\ werden ausgeschlossen.
    \nAbsatz
    \textbf{Forschungsbereich:} Studien aus geisteswissenschaftlichen, sozialwissenschaftlichen oder \\ anderen nicht-naturwissenschaftlichen Disziplinen werden überwiegend ausgeschlossen.
    \nAbsatz
    \textbf{Inhaltsrelevanz:} Publikationen, die keinen direkten Bezug zur Forschungsfrage haben oder thematisch zu weit entfernt sind, werden ausgeschlossen.
    \nAbsatz
    \textbf{Publikationsart:} Populärwissenschaftliche Artikel, Blogbeiträge oder rein kommerzielle \\ Produktbeschreibungen werden nicht berücksichtigt.
    \nAbsatz
\subsubsection{Auswahlprozess der Literatur:} 
\\
Der Auswahlprozess der relevanten Publikationen erfolgt in mehreren Stufen, um eine hohe Präzision zu gewährleisten und den Arbeitsaufwand effizient zu gestalten:
\\
    Screening nach Titel und Abstract: Zunächst werden alle durch die Suchstrategie identifizierten Publikationen anhand ihres Titels und Abstracts begutachtet. Hierbei wird eine erste \\ grobe Einschätzung der Relevanz vorgenommen. Publikationen, die offensichtlich nicht relevant sind, werden an dieser Stelle bereits ausgeschlossen.
\\
    Screening nach relevanten Kapiteln: Bei Titeln und Abstracts, die vielversprechend erscheinen, wird ein tiefergehendes Screening relevanter Kapitel durchgeführt. Dies können Einleitungen, Methodenteile, Diskussionsabschnitte oder Schlussfolgerungen sein, die direkte Hinweise auf die Beantwortung der Forschungsfrage geben könnten.
\\
    Volltextanalyse: Nur Publikationen, die die vorhergehenden Screening-Phasen erfolgreich \\ durchlaufen haben, werden im Volltext analysiert. Dies ist die intensivste Phase, \\ in der die Inhalte detailliert auf ihre Relevanz, Methodik, Ergebnisse und Beiträge zur \\ Forschungsfrage geprüft werden.
\nAbsatz
\textbf{Kriterien für die Priorisierung der Quellen:}
\nAbsatz
    \textbf{Zitationen:} Bei der Auswahl der Quellen wird ein besonderes Augenmerk auf die Anzahl \\ der Zitationen gelegt. Hoch zitierte Arbeiten gelten in der Regel als ein Indikator für \\ wissenschaftliche Qualität und Einfluss innerhalb der Forschungsgemeinschaft. Primär werden diese Quellen begutachtet, da sie oft grundlegende Konzepte oder wegweisende Studien \\ darstellen.
\nAbsatz
    \textbf{Publikationsjahr:} Das Publikationsjahr spielt eine differenzierte Rolle bei der Auswahl. \nAbsatz
       \textbf{ Fallbeispiele und Anwendungsfälle:} Bei Fallbeispielen, aktuellen Anwendungsfällen oder \\ Studien, die sich mit der Leistungsfähigkeit spezifischer, sich schnell entwickelnder Technologien befassen, wird der Zeitraum der letzten 5-10 Jahre als relevanter erachtet. Dies stellt sicher, dass die betrachteten Beispiele dem aktuellen Stand der Technik entsprechen.\nAbsatz
        \textbf{Grundlagen und Definitionen:} Sollte die Suche jedoch nach Definitionen gängiger \\ Methoden, grundlegenden Prinzipien oder etablierten Konzepten erfolgen, deren Kerninhalt über die Zeit stabil geblieben ist, werden auch ältere Texte der letzten 20 Jahre in Betracht gezogen. Hier ist die zeitliche Aktualität weniger kritisch, solange die dargestellten Grundlagen noch valide sind.
\nAbsatz
Durch diese gestufte Vorgehensweise und die differenzierten Auswahlkriterien wird \\ sichergestellt, dass die systematische Literaturrecherche sowohl aktuelle Entwicklungen als auch bewährte Grundlagen berücksichtigt, um eine umfassende und fundierte Beantwortung der \\ Forschungsfrage zu ermöglichen.


 